\scanpage{149}
Compléments aux notes IHES

\deleted{Exercice}

\uline{Remarque} La méthode de démonstration de 1.7. prouve plus généralement l'énoncé suivant : Soit $X$ et $Y$ sur $S$, $X$ propre, \inserted{surjectif,} séparable à fibres univ\textsuperscript{t} connexes, et $Y$ connexe, \inserted{les préschémas loc. \unsure{noeth.}} \unsure{Se donnant|Donnant|Donnons|Désignons} un p\textsuperscript{t} géométrique $c$ de $X\times_S Y$, d'où des p\textsuperscript{ts} géom $a,b,d$ dans $X,Y,S$ et un hom canonique $\pi_1(X\times_S Y,c)\longrightarrow\pi_1(X,a)\times_{\pi_1(S,d)}\pi_1(Y,b)$.

Cet hom est surjectif. \deleted{Si} Il est bijectif si \inserted{$Y$ a une section sur $S$, plus généralement si} la condition suivante est vérifiée : Il existe un $S'$ \deleted{\unsure{étale}} \inserted{connexe,} non vide, \deleted{\unsure{étale}} sur $S$, tel que $Y\times_S S'=Y'$ ait une section sur $S'$ [i.e. \deleted{\unsure{il y a|il faut}} il existe un $S$-morphisme $\unsure{S'|X'}\longrightarrow Y$] \deleted{et que $\pi_1$ fasse une \unsure{correspondance} \unsure{biunivoque}} et que $\pi_1(X')\unsure{-|\to}\pi_1(X)$ soit injectif [par exemple \unsure{fini} $S'$ étale sur $S$]\deleted{\unsure{lorsque}}. Cet énoncé s'applique en particulier chaque fois que $S$ est le spectre d'un corps $k$, et donnera

% slanted
\margin{$S'\longrightarrow S$ tel que\\ \deleted{\unsure{$f$ $X'$ \illegible{2}}}}

\[
\pi_1(X\times_k Y)\xrightarrow{\sim}\pi_1(X)\times_{\pi_1(k)}\pi_1(Y)
\]

[$X$ propre univ\textsuperscript{t} connexe sur $k$, $Y$ \deleted{\unsure{univ\textsuperscript{t}}} connexe sur $k$ sans plus]

% archive sheet number: 148

